\section{Políticas jerárquicas Goal-Conditioned}

\subsection{Subobjetivos de alto nivel y ejecución de bajo nivel}

La política de alto nivel produce subobjetivos continuos o discretos \(g\), y la política de bajo nivel toma \((s, g)\) como conditioning para muestrear acciones. En el continuous case, el subobjetivo puede definirse como un embedding de alta dimensión de los elementos; en el discrete case, puede usarse directamente como el id de una habilidad.

\begin{table}[H]
\centering
\begin{tabular}{p{3cm}p{5cm}p{4cm}}
\toprule
Método & Mecanismo central & Tratamiento típico \\
\midrule
HIRO & Off-policy correction & El nivel alto restablece el objetivo cada $c$ pasos \\
HAC & Recursive hindsight relabeling & La reward del nivel alto proviene directamente de los subobjetivos \\
FuN & Predictive subgoal representations & Usa un predictive model para medir si el subgoal es alcanzable \\
\bottomrule
\end{tabular}
\caption{Comparación de métodos jerárquicos goal-conditioned comunes.}
\end{table}

\begin{knowledgebox}{Consideraciones para diseñar subobjetivos}
\begin{itemize}
    \item Alcanzabilidad: el subobjetivo debe poder lograrse dentro del horizon del nivel bajo;
    \item Identificabilidad: la política de bajo nivel necesita saber cuándo ha completado el subobjetivo;
    \item Reutilización: conviene diseñar subobjetivos generales para evitar que el nivel alto cree objetivos nuevos con frecuencia;
    \item Estabilidad: la distribución de subobjetivos de alto nivel debe ser acotada y se recurre al hindsight relabel para controlar el drift.
\end{itemize}
\end{knowledgebox}

\subsection{Retos del entrenamiento y ciclos de retroalimentación}

El entrenamiento del RL jerárquico suele venir acompañado de non-stationarity: la política de alto nivel se enfrenta a un comportamiento de bajo nivel que cambia constantemente, y la política de bajo nivel depende, a su vez, de subobjetivos dinámicos. El ponente recomienda usar un replay buffer de doble vía (uno para high y otro para low) y escribir un log que registre el hash del high-level goal y la associated reward.

\begin{warningbox}{Propagación del gradiente y conflictos en la repetición de experiencias}
\begin{itemize}
    \item el low-level replay puede sobreajustarse a los high-level plans tempranos;
    \item la high-level reward requiere un hindsight relabel basado en la low-level completion; de lo contrario, el ruido del gradiente es demasiado grande;
    \item usar off-policy correction (como en HIRO) permite corregir el bias del high-level update.
\end{itemize}
\end{warningbox}

\subsection{Hindsight relabeling y reutilización de datos}

El hindsight relabeling toma los subobjetivos no alcanzados del high-level plan y los «reescribe» como un nuevo pseudo-goal a partir de los estados que el nivel bajo sí alcanzó, de modo que el replay buffer aporte más muestras positivas. Al implementarlo, hay que controlar la frecuencia mediante un relabel scheduler para evitar que un rewrite excesivo dañe el original objective; los metadata de todas las relabeled samples deben escribirse en el prompt log para permitir su audit.

\begin{importantbox}{Pasos para implementar el hindsight relabel}
\begin{enumerate}
    \item Recolectar el goal de alto nivel y la trajectory de bajo nivel de un episode;
    \item Elegir un estado como relabel goal;
    \item Recalcular la reward y actualizar la Q de alto nivel;
    \item Escribir los metadata de la relabeled sample (goal hash + timestamp) en la instrumentation;
\end{enumerate}
\end{importantbox}

\subsection{Resumen de la sección}

Las políticas jerárquicas goal-conditioned hacen que el nivel alto fije subobjetivos como un comandante y tratan al nivel bajo como el ejecutor; la dificultad del entrenamiento reside en el drift frecuente del entorno y de los objetivos, lo que exige guardrails de hindsight y off-policy.
